\chapter{Introduction}
[\textit{Must be present in FYDP-1 Report and also in Final Report}]


This chapter introduces the research problem, explains why it is important in a real computational setting, and defines the objectives and expected outcomes of the project. It also presents a concise summary of the proposed methodology and outlines the structure of the remaining report.

\section{Project Overview}
Ordinary differential equations (ODEs) are used to model time-dependent phenomena across physics, engineering, and applied mathematics. Many physically meaningful systems, including atmospheric circulation, fluid interaction, and coupled dynamical processes, are naturally expressed as initial value problems whose states evolve continuously with time. In practice, such systems are usually solved with numerical techniques such as Euler-type methods and Runge-Kutta schemes because closed-form solutions are rarely available for nonlinear coupled equations.

This project focuses on the Lorenz-1960 ODE system, a reduced meteorological model that represents nonlinear interactions among three coupled state variables. The system is a useful research case because it is small enough to study carefully, but still complex enough to expose the strengths and weaknesses of different approximation strategies. Existing numerical solvers can produce accurate trajectories for the Lorenz-1960 system, yet they must be executed again for every new setup and do not directly provide a reusable learned mapping from time to state values. Recent machine learning research has shown that neural networks can approximate the solutions of differential equations, either by embedding physical laws into the loss function or by learning solution functions directly from data \cite{matthews2025pinnde,pratama2022ann,chen2019neural}.

The problem addressed in this project is therefore not whether the Lorenz-1960 system can be solved at all, but whether a plain artificial neural network (ANN) can learn this solution accurately enough to serve as a practical alternative approximation model, and which ANN architecture performs best for this task. The study is framed around architecture selection, because current related work does not clearly establish how depth, width, and activation choice affect plain ANN performance on this specific ODE system. The proposed work therefore aims to identify an ANN architecture that approximates the Lorenz-1960 solution with high accuracy while remaining simpler and easier to train than physics-informed alternatives.

\section{Motivation}
The computational motivation of this work comes from the growing need for efficient surrogate models for scientific computation. Classical numerical methods are highly reliable, but they are solver-driven rather than model-driven: when the equation, the time interval, or the initial condition changes, the numerical process must be repeated. This is acceptable for one simulation, but it becomes expensive in repeated experimentation, parameter sweeps, and learning-based pipelines where many evaluations may be required. A trained ANN offers a different possibility by learning a direct approximation from time to solution state, allowing rapid inference after training.

There is also a clear research motivation. In the recent literature, one branch of work emphasizes physics-informed neural networks and operator-learning frameworks, while another branch studies pure ANN-based approximators for differential equations. Physics-informed approaches are attractive because they inject equation residuals into the learning objective, but they also introduce additional optimization complexity and hyperparameter balancing requirements. Since the supervisor guidance already points to PINNs as an important reference direction, a meaningful and well-scoped FYDP contribution is to investigate the opposite design choice: removing the physics term and evaluating how far a plain ANN can go on a well-defined ODE benchmark. This makes the project both educational and scientifically useful because it clarifies whether model architecture alone can deliver competitive approximation quality on the Lorenz-1960 problem.

The project is also beneficial from an application perspective. Understanding how ANN architectures behave on a nonlinear physical system can support later work on more complex scientific machine learning tasks, including surrogate modeling, reduced-order simulation, and data-driven forecasting. Even if the final conclusion is that plain ANNs are insufficient for this class of problem, that result is still valuable because it helps justify when more advanced physics-informed or operator-learning models are actually necessary.

\section{Objectives}
The objectives of this study are fivefold. First, the project aims to formulate the Lorenz-1960 ODE system as a supervised approximation problem by generating a high-accuracy numerical reference solution over the target time interval. Second, it aims to design and train multiple ANN architectures that take time as input and predict the corresponding state variables of the system. Third, it seeks to compare these architectures systematically by varying depth, width, activation function, and training strategy in order to measure their effect on approximation accuracy. Fourth, it aims to identify the architecture that achieves the best overall performance against the numerical reference solution. Fifth, it intends to position the resulting best architecture against the most relevant published approaches, especially the physics-informed DeepONet-based treatment of the same Lorenz-1960 system and the broader ANN literature on differential equation solving.

\section{Methodology}
The methodology of the project is based on a comparative experimental design. The Lorenz-1960 system will first be defined as an initial value problem using the parameter setting discussed in the reference study. A high-accuracy numerical solution will then be generated using a Runge-Kutta based solver and treated as ground truth for model training and evaluation. After that, a family of feedforward ANN models will be constructed with different numbers of hidden layers, different neuron counts, and different activation functions. Each model will be trained to learn the mapping from time to the three solution components of the Lorenz-1960 system.

Once training is complete, the candidate architectures will be evaluated with quantitative error metrics such as mean squared error, along with convergence behavior, parameter count, and inference efficiency. The best-performing configuration will then be analyzed in relation to the recent literature in order to understand whether a plain ANN can provide a useful approximation baseline for this problem and how its strengths and limitations compare with more structured neural differential equation methods.

\section{Project Outcome}
The expected outcome of this project is an evidence-based identification of an ANN architecture that is most suitable for approximating the Lorenz-1960 ODE solution within the chosen experimental setting. More specifically, the project should produce a reproducible dataset generation process, a set of trained ANN baselines, a comparative analysis of architectural choices, and a final recommendation regarding which configuration gives the best trade-off between simplicity and accuracy.

In addition to the best-model result, the project is expected to produce a clearer scientific positioning of plain ANN solvers within the broader landscape of differential equation learning methods. If the proposed ANN models perform well, the work will demonstrate that a relatively simple data-driven solver can serve as a useful approximation tool for this benchmark system. If the performance is weak or unstable, the work will still yield a valuable conclusion by showing why additional structure, such as physics-informed losses or operator-learning designs, may be necessary. In either case, the project outcome will be a justified research contribution rather than only an implementation exercise.

\section{Organization of the Report}
The remainder of this report is organized as follows. Chapter 2 presents the theoretical background required for the study, including ODE fundamentals, ANN concepts, and the most relevant related research on neural approaches for differential equations. Chapter 3 will describe the project design in detail, including requirement analysis, the selected solution strategy, project planning, and task allocation. Chapter 4 will present the implementation of the proposed system and the experimental setup used to train and evaluate the candidate architectures. Chapter 5 will discuss the obtained results, compare the architectures, and interpret the findings in relation to the stated objectives and the literature. Finally, Chapter 6 will summarize the overall contribution, highlight the main limitations, and suggest directions for future work.
